\chapter{Propensity Score Matching}
\label{cap:psm}
% ── index ──────────────────────────────────────────────────────
\index{propensity score matching|textbf}
\index{PSM|see{propensity score matching}}
\index{psm@\texttt{psm()}|textbf}
\index{propensity score|textbf}
\index{ATT|see{average treatment effect on the treated}}
\index{caliper}
\index{selection bias}

% ─────────────────────────────────────────────────────────────────────────────
\section{Selection on Observables}

In many observational studies, treatment assignment is not random
— individuals select into treatment based on observable characteristics
$X$. If $X$ also affects the outcome $Y$, direct comparison between
treated and controls produces biased estimates.

\textbf{Propensity Score Matching} (PSM) exploits the Rosenbaum and
Rubin (1983) theorem: if treatment assignment is independent of potential
outcomes given $X$ (\emph{ignorability}), then the same independence holds
given the propensity score $p(X) = P(D=1 \mid X)$:

\[
  (Y^0, Y^1) \perp D \mid X
  \quad\Longrightarrow\quad
  (Y^0, Y^1) \perp D \mid p(X)
\]

This reduces the high-dimensional problem (conditioning on $X$) to a
one-dimensional problem (conditioning on $\hat{p}$). PSM:
\begin{enumerate}
  \item Estimates $\hat{p}(X_i)$ via logit (propensity score).
  \item Matches each treated unit with the closest control(s) in $\hat{p}$.
  \item Computes the ATT as the mean difference in the matched sample.
\end{enumerate}

The common support (\emph{overlap}) condition: $0 < p(X) < 1$ for all $X$
is essential — units with probability 0 or 1 of treatment have no counterfactual
in the other group.

% ─────────────────────────────────────────────────────────────────────────────
\section{DGP Verification}

The DGP creates confounding via $x_1$: those with high $x_1$ have a higher
probability of treatment and also a higher outcome on average. The true ATT is $\tau = 2$.

\[
  P(D=1 \mid x_1) = 0{,}3 + 0{,}4\,x_1 \in [0{,}30,\; 0{,}70]
  \qquad
  y_i = 1 + 2\,d_i + 3\,x_{1i} + \varepsilon_i
\]

\begin{lstlisting}[caption={DGP with selection on observables}]
seed(42)
input df
id
1
end
for i in 1..=10 { df = append(df, df) }
df |> filter(_n <= 1000)
generate df x1 = runiform()                   // x1 ~ U(0,1)
generate df u  = runiform()                   // assignment noise
generate df d  = (0.3 + 0.4*x1) > u          // P(D=1|x1) ∈ [0.30, 0.70]
generate df e  = rnormal()
generate df y  = 1.0 + 2.0*d + 3.0*x1 + e   // true ATT = 2
\end{lstlisting}

\subsection*{Naive OLS vs.\ OLS with controls}

\begin{lstlisting}[caption={OLS comparison without and with x1 control}]
let m_ols = ols(y ~ d, df)          // no control: biased
let m_adj = ols(y ~ d + x1, df)    // with control: unbiased
display m_ols
display m_adj
\end{lstlisting}

\begin{lstlisting}[language={}, basicstyle=\ttfamily\small,
  frame=none, numbers=none, backgroundcolor=\color{codbg},
  xleftmargin=0pt, xrightmargin=0pt, breaklines=false]
── OLS without control ───────────────────────────────
const  2.9135***  (0.0398)    d  2.2215***  (0.0562)
N=1000  R²=0.6105
──────────────────────────────────────────────────────

── OLS with x1 ───────────────────────────────────────
const  1.5289***  (0.0193)    d  1.9762***  (0.0180)
x1     3.0053***  (0.0317)
N=1000  R²=0.9610
──────────────────────────────────────────────────────
\end{lstlisting}

OLS without control overestimates $\tau$ by $+0{,}22$ (2.22 vs.\ 2.0) — omitted
variable bias. Controlling for $x_1$ the OLS recovers $\hat\tau = 1{,}98$,
indistinguishable from 2.

\subsection*{PSM without caliper}

\begin{lstlisting}[caption={1:1 PSM without caliper}]
let m_psm = psm(y ~ d + x1, df, k=1, boot=200)
display m_psm
\end{lstlisting}

Without a caliper, matching accepts pairs that are distant in $\hat{p}$. Balance
may not improve much and the estimate can remain biased if treated units are
matched to poor controls.

\subsection*{PSM with caliper}

The caliper limits the maximum distance allowed in matching. The practical rule
of Cochran and Rubin (1973) is caliper $= 0{,}2 \times \sigma_{\hat{p}}$, which
corresponds approximately to $0{,}05$ for uniform scores.

\begin{lstlisting}[caption={PSM with caliper = 0.05}]
let m_cal = psm(y ~ d + x1, df, k=1, caliper=0.05, boot=200)
display m_cal
\end{lstlisting}

With caliper $= 0{,}05$, the estimate should move closer to the true ATT at
the cost of dropping treated units with no acceptable match. The caliper should
also improve balance by rejecting distant propensity-score matches.\footnote{%
  The \texttt{SMD > 0.10} warning persists because the conventional rule requires
  SMD $\leq 0{,}10$ (Austin 2009). In this DGP the confounding by $x_1$ is intense
  ($\beta_{x_1} = 3$) and $N = 1000$; with more observations the caliper could
  be narrowed to further improve balance.}

% ─────────────────────────────────────────────────────────────────────────────
\section{Syntax}

\begin{lstlisting}
psm(outcome ~ treatment + cov1 + cov2, df)
psm(outcome ~ treatment + cov1 + cov2, df, k=1, caliper=0.05, replace=false, boot=200)
\end{lstlisting}

The formula: \lstinline{outcome ~ treatment + cov1 + cov2 + ...}. The first
regressor after \lstinline{~} is the treatment variable; the rest are the
covariates used in the propensity model.

\begin{center}
\begin{tabular}{lll}
\toprule
\textbf{Option} & \textbf{Default} & \textbf{Description} \\
\midrule
\texttt{k}        & 1       & number of controls per treated unit \\
\texttt{caliper}  & none    & maximum distance in propensity score \\
\texttt{replace}  & false   & matching with replacement \\
\texttt{boot}     & 200     & bootstrap replicas for SE \\
\bottomrule
\end{tabular}
\end{center}

\begin{lstlisting}[caption={Typical workflow — program evaluation}]
load "beneficiarios.csv" as df

// Selection covariates
let m = psm(renda ~ tratado + idade + educ + regiao, df,
            k=1, caliper=0.05, boot=500)
display m
\end{lstlisting}

\begin{nota}
  PSM is consistent under \emph{ignorability} (selection on observables):
  $\{Y^0, Y^1\} \perp D \mid X$. Omitted variables that simultaneously affect
  $D$ and $Y$ invalidate identification — the same problem that makes IV
  necessary in ch.~\ref{cap:iv}. Unlike DiD
  (ch.~\ref{cap:did}), PSM does not eliminate time-constant unobserved confounders.
\end{nota}

\section{Empirical validation}

The Wooldridge \texttt{jtrain3} dataset is used in the \texttt{psm\_jtrain3}
validation case in \texttt{validation/cases/}. The case estimates the ATT of
a job-training programme by PSM in \hayashi{} and compares it with a Python
reference implementation. Run \lstinline{hay validate} to see the current
status.
